\documentclass[12pt]{article}
\usepackage[a4paper,margin=2.2cm]{geometry}
\usepackage{amsmath,amssymb,amsthm,mathtools}
\usepackage{hyperref}
\usepackage{cleveref}

% 中文（xelatex）
\usepackage{fontspec}
\usepackage{xeCJK}
\setCJKmainfont{Noto Serif CJK SC}

\title{Einstein Sum}
\date{}
\begin{document}
\maketitle

\vspace{0.5em}
\noindent\textbf{1: 向量的点乘（1维）}\\
向量 $u,v \in \mathbb{R}^{n}$,
使用 $\sum$ 写出 $\vec{u} \cdot \vec{v}$ 的点积式

$\vec{u} \cdot \vec{v}$  = $\sum_{i = 1}^{n}u_{i}v_{i} = u_{i}v_{i}(Einstein Sum)$
i 是 哑标 不存在 自由标


\vspace{0.5em}
\noindent\textbf{2: 矩阵乘向量}\\
矩阵 $A \in \mathbb{R}^{m \times n}$,$\vec{v} \in \mathbb{R}^{n}$ \\
$\vec{y} = Ax, \;\; \vec{y} \in \mathbb{R}^{m}$

1: 使用 $\sum$ 符号写 $y_{i}$

$y_{i} = \sum_{j = 1}^{n}A_{ij}x_{j}$


2: 使用 爱因斯坦求和约定

$y_{i} = A_{ij}x_{j}$

3: 哪个字母是自由指标？哪个字母是哑标?

i 是自由标, j 是哑标

\vspace{0.5em}

向量 $\vec{v} \in \mathbb{R}^{1 \times m}$(也是矩阵) \\
$B \in \mathbb{R}^{m \times n}$ \\
$\vec{z} = vB, \;\; \vec{z} \in  \mathbb{R}^{1 \times n}$ \\

1: 使用 $\sum$ 符号写 $z_{k}$

$z_{k} = \sum_{j = 1}^{m}v_{1j}B_{jk}$


2: 使用 爱因斯坦求和约定
$z_{k} = v_{1j}B_{jk}$


3: 哪个字母是自由指标？哪个字母是哑标?

k 是自由标, j 是哑标


\vspace{0.5em}
\noindent\textbf{3: 矩阵乘矩阵}

$A \in \mathbb{R}^{m \times p}$, $B \in \mathbb{R}^{p \times n}$
$C = AB, C \in \mathbb{R}^{m \times n}$

矩阵 C 的 第 i 行, 第 k 列, 记 $C_{ik}$

1:使用 $\sum$ 符号写 $C_{ik}$

$C_{ik} = \sum_{j = 1}^{p} A_{ij}B_{jk}$

2: 使用 爱因斯坦求和约定
$C_{ik} = A_{ij}B_{jk}$

3: 哪个字母是自由指标？哪个字母是哑标?

ik 是自由标, j 是哑标


\vspace{0.5em}
\noindent\textbf{5: tr(XY) = tr(YX)}

$X,Y \in \mathbb{R}^{n \times n}$

1: $tr(XY)$

\begin{align}
    XY_{ik} = \sum_{j = 1}^{n}X_{ij}Y_{jk}                               \\
    tr(XY) =  \sum_{j = 1}^{n}XY_{ii}                                    \\
    tr(XY) = \sum_{i = 1}^{n}(\sum_{j = 1}^{n}X_{ij}Y_{ji}) = \mathbb{R} \\
    tr(XY) = X_{ij}Y_{ji}(EinsteinSum)
\end{align}




\end{document}